\documentclass[landscape,twocolumn,leqno,a4paper,10pt]{report} %leqno affiche les numéros d'équations à gauche

\usepackage[a4paper,top=0.4cm,bottom=1cm,left=0.4cm,right=0.4cm,marginparwidth=0.5cm,footskip=0.6cm]{geometry}



%----------------- Codage de la langue et de la police--------------------------------------------------------------------------
\usepackage[french]{babel} 		% Permet de spécifier au compilateur que l'on désire taper son document en français.
\usepackage[utf8x]{inputenc} 	% Permet de taper directement à l'écran, dans le fichier .tex, les caractères accentués.
\usepackage[T1]{fontenc} 		% Permet juste d'utiliser une nouvelle norme LATEX concernant le codage des caractères.
%\usepackage{lmodern}		% Pack de police --> Pas besoin puisque j'utilise Fourier

%------------------------ Packages utiles-------------------------------------------------------------------------------------------

\usepackage{graphicx}	% Pour insérer des images
\usepackage{caption}		% Légende

\newcommand\hmmax{0}		% Alors c'est une astuce assez mystérieuse pour dépasser la limite de polices latex.
							% En effet, le package bm est lourd et avec tout ce que l'on met, on obtient un message d'erreur : Too many math alphabets used in version normal.			
\newcommand\bmmax{0}		% Cette astuce règle le problème, il faut tout de même mettre ces compteurs avant le package bm

\usepackage{amsmath,amssymb}	% Permettent d'obtenir des caractères mathématiques (de toute beauté !).
\usepackage{amsthm}		% Pour pouvoir faire les newthm
\usepackage{mathtools}	% Corrige des bugs dans amsmath et le complète
\usepackage{numprint}	% Pour les espacements dans les nombres : \numprint{10000} = 10 000
\usepackage{fourier}		% REMARQUE : [upright] impose aux majuscules d'être droite dans le mode mathématique 
\usepackage{pifont} 		% Importe une large gamme de symboles, on utilise \ding{numéro}
\usepackage{xargs} 		% Permet de définir des macros à plusieurs paramètres optionnels.
\usepackage{fancyhdr} 	% Package pour les entêtes et bas de page
\pagestyle{fancy} 			% Accompagne \usepackage{fancyhdr}
\usepackage{bm} 		% Permet de mettre en gras du contenu mathématiques via \bm{..}
\usepackage{mathrsfs} 	% Pour les lettres cursives
\usepackage{enumitem} 	% Pour pouvoir personnaliser les listes

\frenchbsetup{StandardLists=true} % Afin d'éviter tout conflit entre les extensions enumitem et french babel. 

\usepackage{sectsty}	% Pour colorer les titres
\usepackage{titlesec} 	% Pour personnaliser les titres
\usepackage{setspace} 	% Pour changer la taille des interlignes
       	               
\usepackage{array,multirow,makecell}
\usepackage{hlist}		% Package Christian Tellechea sur les listes numérotées gauche vers droite
\setdefaulthlist{left margin = 2mm, label sep= 0.2mm}	% réglages package hlist
				

%------------ Les packages de dessin sous pstricks : il faut changer la compilation de latex pour compiler ces dessins------------------
% Options\Configurer Texmaker\Compil rapide\ Latex+dvips+ps2pdf+Voir PDF

\usepackage[pdf]{pstricks} 	% Pour faire des dessins
\usepackage{pstricks-add,pst-plot,pst-tree,pst-node,pst-math,pst-eucl,pst-func,pst-3dplot}% ,pst-text,pst-eps,pst-fill,pst-grad
\usepackage{multido}		% Pour faire des multido
\usepackage{lastpage} 	% Pour le compteur de la dernière page, on appelle alors celle-ci avec lastage



\usepackage[colorlinks=true, 
			%allcolors=blue
			]{hyperref}		% ATTENTION : doit être importé après la majorité des packages sauf tabularx et breakrules (par exemple)
\hypersetup{			% Paramétrage des hyperliens
	colorlinks = true,   		% Colorise les liens
	breaklinks = true,		% Permet les retours à la ligne pour les liens trop longs
	urlcolor = blue,		% Couleur des hyperliens
	linkcolor = black,		% Couleur des liens internes aux documents (index, figures, tableaux, équations, ...)
	citecolor = blue		% Couleur des liens vers les références bibliographiques
	}




%%%%%%%%%%%%%%%%%%%%%%%%%%%%%%%%%%%%%%%%%%%%%%%%%%
%------------------------- Définition des commandes -----------------------------------------------------------------------------
%%%%%%%%%%%%%%%%%%%%%%%%%%%%%%%%%%%%%%%%%%%%%%%%%%


%--------------------- Raccourcis fontes--------------------------------------------------------------
% REMARQUE : le package fourier met la même police pour mathcal et mathscr, on remet donc mathcal normalement et on utilise la définition de fourier (un peu plus grasse) pour mathscr
\DeclareMathAlphabet{\mathcal}{OMS}{cmsy}{m}{n}	
\DeclareMathAlphabet{\mathscr}{FMS}{futm}{m}{n}

%---------------------------- Raccourcis des espaces--------------------------------------------------------------------------------------------
\newcommand\n{\noindent}
\newcommand\q{\quad}
\newcommand\qq{\quad\quad}

%--------------------- Raccourcis en mode mathématiques-------------------------------------------------------------------------------------
\newcommand{\et}{\text{ et }}
\newcommand{\ou}{\text{ ou }}
\newcommand{\SSI}{\Longleftrightarrow}
\let\si\relax
\newcommand{\si}{\text{ si }}
\newcommand{\sinon}{\text{ sinon }}
\renewcommand{\epsilon}{\varepsilon}		% epsilon beaucoup plus joli
\renewcommand{\geq}{\geqslant} 
\renewcommand{\leq}{\leqslant}
\newcommand{\e}{\mathrm{e}}
\renewcommand{\textbf}[1]{\begingroup\bfseries\mathversion{bold}#1\endgroup}		% Très important, permet de passer en gras même avec le mode mathématiques

\newcommand{\Mod}[1]{\allowbreak\if@display\mkern5mu\else\mkern5mu\fi\left[#1\right]}	% Pour modulo entre crochets [.] avec un espacement convenable 


%---------------------------------- Raccourcis calcul diff----------------------------------------------------------------------------------------
\renewcommand{\d}{\;\mathrm{d}}	% Pour le d droit de l'intégrale car ce n'est pas une variable...


%-------------------------------- Raccourcis des ensembles----------------------------------------------------------------------------
\newcommand\R{\mathbb{R}}
\newcommand\C{\mathbb{C}}
\newcommand\K{\mathbb{K}}
\newcommand\Q{\mathbb{Q}}
\newcommand\D{\mathbb{D}}
\newcommand\Z{\mathbb{Z}}
\newcommand\N{\mathbb{N}}
\newcommand\U{\mathbb{U}}


%------------------------------- Dimensions intelignes, items...-----------------------------------------------------------------------------------------------
\onehalfspacing	% Pour que l'interligne soit de 1.5. Pour que l'interligne soit de la valeur que l'on souhaite, de base : \setstretch{1}

\setlist[1]{font=\bfseries,itemsep=0.2mm,leftmargin=6mm}
\setlist[2]{font=\bfseries,itemsep=0.2mm,leftmargin=4mm}
\setenumerate[1]{font=\bfseries,itemsep=0.2mm,leftmargin=6mm}
\setenumerate[2]{font=\bfseries, label=\alph*., leftmargin=4mm}
%%\setenumerate[profondeur]{font=\bfseries, label=\ding{221} \color{blue}} permet de changer en a) b) les sous-items d'une liste enumerate
%% itemsep = correspond à l'interligne entre chaque item
% de base : labelsep=2mm et permet de changer la distance entre l’étiquette et le texte des éléments
% leftmargin =4mm spécifie la marge à gauche de la liste
\AtBeginDocument{\renewcommand{\labelitemi}{\textbullet}}	% Pour que les puces par défaut soit des points.
%\AtBeginDocument{\renewcommand{\labelitemi}{\ding{86}}}	% Pour que les puces par défaut soit des étoiles.
\AtBeginDocument{\renewcommand{\labelitemii}{$\ast$}}		% Pour que les puces par défaut pour une sous-liste soit des astérisques.
\newcommand{\itemmm}{\item[$\ast$]}				% \itemmm devient l'item \ast
\newcommand{\itemmmm}{\item[\footnotesize$\circledcirc$\normalsize]}
\newcommand{\itemmmmm}{\item[\footnotesize$\astrosun$\normalsize]}
\newcommand{\itemmmmmm}{\item[\footnotesize$\multimap$\normalsize]}
\newcommand{\itemmmmmmmmm}{\item[$\divideontimes$]}
\newcommand{\etoile}{\item[$\star$]}
\newcommand{\etoilee}{\item[$\bigstar$]}
\newcommand{\losange}{\item[$\diamondsuit$]}
\newcommand{\carre}{\item[\LARGE$\centerdot$]}
\newcommand{\carreombre}{\item[\ding{112}]}

\setlength{\columnsep}{30pt} 		% Permet d'écarter les deux colonnes du twocolumn de 30pt, pratique pour écarter le texte de la ligne centrale
\setlength{\columnseprule}{0.5pt}	% Affiche une ligne verticale entre chaque item et entre chaque page du twocolumn
						% \columnseprule=0cm 	% Permet d'enlever le trait vertical entre chaque colonne
%\renewcommand\tabenumsep{0.5cm\hskiplem}	% Ecarte de 0.5cm les items de tabenum, sinon ils sont trop collés...


\newlength{\tagmarginsep} 		% On crée une longueur
\setlength{\tagmarginsep}{0cm}		% Permet de régler la position d'un tag pour une équation mettre avant l'équation : \setlength{\tagmarginsep}{x cm} pour que le tag se rapproche de l'équation
\everydisplay{\displayindent=\tagmarginsep \displaywidth=\dimexpr\linewidth-2\tagmarginsep}	% Va avec ce qu'il y a juste avant


%-------------------- Définition des couleurs-----------------------------------------------------------------------------------------

\definecolor{violetfonce}{RGB}{126,62,142}
\definecolor{magentafonce}{RGB}{128, 0, 128}
\newrgbcolor{magentafonce}{.5 0 .5}			% Une autre façon de définir une couleur, permet surtout d'utiliser \magenta
\definecolor{violet1}{cmyk}{0.5,1,0,0}
\definecolor{violet2}{cmyk}{0.24998, 1, 0, 0}
\definecolor{violet3}{cmyk}{0, 1, 0.25, 0}



\newcommand\tred{\textcolor{red}}
\newcommand\tredf{\textcolor{rougefonce}}
\newcommand\tdred{\textcolor{deepred}}
\newcommand\tblue{\textcolor{blue}}
\newcommand\tbluef{\textcolor{bleufonce}}
\newcommand\tgreen{\textcolor{vertfonce}}
\newcommand\torange{\textcolor{orange}}
\newcommand\tmagenta{\textcolor{magentafonce}}


%------------------ Changement paramètres des parties, sous parties---------------------------------------------------------


\sectionfont{\color{violet1}}			% Utilisation du package sectsty
\subsectionfont{\color{violet2}}
\subsubsectionfont{\color{violet3}}


%--------------------------------- Raccourcis probabilités--------------------------------------------------------------------------------
\newcommand{\Prob}{\mathbb{P}}
\newcommand{\E}{\mathbb{E}}
\newcommand{\normale}[2]{\mathcal{N}\left(#1;#2 \right)}	% Pour la loi Normale : 	$\normale{\mu}{\sigma^2}$
\newcommand{\Unif}[1]{\mathcal{U}\left(#1\right)}			% Pour la loi uniforme :	$\Unif{[a;b]}$
\newcommand{\Ber}[1]{\mathcal{B}\left(#1\right)}			% Pour la loi de Bernoulli : $\Ber{p}$
\newcommand{\Bino}[2]{\mathcal{B}\left(#1;#2\right)}		% Pour la loi Binomiale :	$\Bino{n}{p}$
\newcommand*{\binomial}[2]{% 
  \ensuremath{ 
    \begin{pmatrix} 
      #2\\ 
      #1
    \end{pmatrix}}}				% Pour le coefficient binomial : $\binom{k}{n}$
\renewcommand{\binom}{\binomial} 	% La commande \binom existait déjà mais n'était pas pratique



%----------------------------------- Opérateurs---------------------------------------------------------------------------------------------------

\newcommand{\barre}[1]{\overline{#1}}  									% Barre au dessus du texte, pratique pour le conjugué des complexes
\newcommandx{\Sum}[3][3=\white]{{\displaystyle\sum\limits_{\substack{#1 #3}}^{#2}}}		% $\Sum{k=0}{n}[<option>]$ où [<option>] est placée à la suite de k=0 si l'on tape de façon classique, et en-dessous si l'on met [\\ k\neq k_0]
															% On peut juste taper : $\Sum{k=0}{n}$ et cela fonctionne
															% \displaystyle impose de gérer les tailles comme en mode mathématique : ici on aura donc toujours de grande somme que ce soit en math comme dans le texte
\newcommandx{\Prod}[3][3=\white]{{\displaystyle\prod\limits_{\substack{#1 #3}}^{#2}}}		% $\Prod{k=0}{n}[<option>]$	idem pour l'option
\newcommand{\Int}[4]{{\displaystyle\int_{#1}^{#2} #3  \d #4}}   					% $\Int{0}{\pi}{\cos(t)}{t}$
\newcommandx{\Lim}[3][3=\white]{\lim\limits_{\substack{#1 \to #2 #3}}}		% $\Lim{x}{0}[<option>]$	idem pour l'option


%------------------------------------ Délimiteurs--------------------------------------------------------------------------------------------------
%\newcommand{\abs}[1]{\left\lvert #1 \right\rvert}			% Valeur absolue | |
%\newcommand{\norm}[1]{\left\lVert #1 \right\rVert}			% Doubles barres verticales : norme d'un vecteur || ||

%--------------------------------------- Vecteurs----------------------------------------------------------------------------------------------------
 
\newcommand{\Oij}{\left(O;\vec{\imath},\vec{\jmath}\right)}
\newcommand{\Oijk}{\left(O;\vec{\imath},\vec{\jmath},\vec{k}\right)}
\newcommand{\Ouv}{\left(O;\vec{u},\vec{v}\right)}
\newcommand{\Ouvw}{\left(O;\vec{u},\vec{v},\vec{w}\right)}

%------------------------------------ Systèmes d'équations-------------------------------------------------------------------------------------------
% Remarque \syst et \syste ne font pas d'alignement mais ils sont bien, si l'on veut de l'alignement, on utilisera \systeme du package "systeme"
\newcommand*{\syst}[2]{%			% Systèmes 2 éq. à ? inconnues : 	\syst{ax+by+c=0}{a'x+b'y+c'=0}
  \left\{\begin{array}{ll} 
        #1\\ 
        #2 
      \end{array}  
    \right. }
 
\newcommand*{\syste}[3]{% 			% Systèmes 3 éq. à ? inconnues : 	\syste{ax+by+c=0}{a'x+b'y+c'=0}{a''x+b''y+c''=0}
  \left\{\begin{array}{ll} 
        #1\\ 
        #2\\
        #3 
      \end{array}  
    \right.}



%------------------------------------------ Autre----------------------------------------------------------------------------------------------------------
\newcommand{\boite}[2]{\noindent\parbox[]{#1 cm}{#2} }  									% Boite autre façon de faire minipage en gros


%---------------------------- En-têtes et bas de pages----------------------------------------------------------------------------
% A chaque fois : *head* pour le haut et *foot* pour le bas
% L = left, C = center, R = right 
% L = toutes les pages, LE = toutes les pages paires, LO = toutes les pages impaires, idem CE, CO, ...


\newcommand{\fcours}[1]{								% f pour fancy				   
	\pagestyle{fancy}									% Efface le style qui a pu être mis avant
	\renewcommand{\headrulewidth}{0pt} 
	\renewcommand{\footrulewidth}{0.5pt} 
	%\fancyhead[L]{}									
	%\fancyhead[C]{}
	%\fancyhead[R]{}									
	\fancyfoot[L]{\vspace{-1mm}\footnotesize \textit{#1}} 			% Classe : <Seconde> 
	\fancyfoot[C]{\vspace{-1mm}\footnotesize \thepage/\pageref{LastPage}}
	\fancyfoot[R]{\vspace{-1.8mm}\textit{\footnotesize C. Marguerite}}	% N'apparaîtra pas si on met les marges avec moins de bottom=1.3cm dans le package geometry
 }												% Fin de la commande
					%% Mettre au début du document : \fcours{Seconde} par exemple
					%% \usepackage[a4paper,top=0.4cm,bottom=1cm,left=0.4cm,right=0.4cm,marginparwidth=0.5cm,footskip=0.6cm]{geometry} 




%---------------------------- Réglages chapitres + table des matières------------------------------------------------------
\titleformat{\part}[display]
{\centering\normalfont\Huge\bfseries}
{\partname}
{20pt}
{\rule{6cm}{1pt}\raisebox{-2.1pt}[10pt][10pt]{\quad\decofourleft\decoone\decofourright\quad}\rule{6cm}{1pt} ~\\\vspace{20pt}}

\makeatletter
\renewcommand\chapter{\global\@topnum\z@
                    \@afterindentfalse
                    \secdef\@chapter\@schapter}
\def\@topnewpage[#1]{#1}
\makeatother										% Pour avoir le chapitre sur une demie page

\titleformat{\chapter}{\centering\normalfont\Large\bfseries\scshape\color{magentafonce}}{}{1em}{}
\titlespacing*{\chapter}{0pt}{-18pt plus 0pt minus 2pt}{0pt plus 2pt minus 2pt}		% Réglages espacements (à gauche, au-dessus, en-dessous) de la commande \chapter

\makeatletter
\newcommand{\twocolumntoc}{		
  \section*{\huge\centering\contentsname
    \@mkboth{%
      \MakeUppercase\contentsname}{\MakeUppercase\contentsname}}%
  \@starttoc{toc}%
}
\makeatother										% Permet de mettre la table des matières en twocolumn 
												% Mettre après \begin{doc..} \twocolumntoc \clearpage 		% Remplace \tableofcontents par \twocolumntoc en twocolumn

\makeatletter
\renewcommand{\thesection}{\@arabic\c@section}
\makeatother										% Enlève la numérotation des chapitres devant les sections

\setcounter{tocdepth}{3} 								% Permet d'afficher dans la table des matières suivant une certaine profondeur : -1=part, 0=chapter, 1=section, ...
\setcounter{secnumdepth}{3}							% Pour avoir la numérotation des subsubsections avec les classes book et report

\newcommand{\nocontentsline}[3]{}
\let\oldaddcontentsline\addcontentsline
\newcommand{\tocless}[2]{\let\addcontentsline\nocontentsline #1{#2}\let\addcontentsline\oldaddcontentsline}	% Pour retirer une section de la table des matières : \tocless\section{...}





\setstretch{1}
\begin{document}
\fcours{Concours général}

\chapter{Concours général des lycées - Session 2021}

\section{Problème 1 : Le début justifie la fin}


\n Dans cet exercice, on considère l'ensemble, noté $\mathscr{S}$, des suites $\left(u_{n}\right)_{n \geq 0}$ à valeurs réelles et telles que 
$$u_{n+1}=\frac{\exp \left(u_{n}\right)}{n+1}$$
pour tout entier $n \geq 0$.\\
Pour tout nombre réel $x$, on note $u(x)$ la suite appartenant à $\mathscr{S}$ et dont le premier terme vaut $x$. On note également $u_{n}(x)$ le terme d'indice $n$ de cette suite. Ainsi, $u_{0}(x)=x$ et $u_{1}(x)=\exp (x)$.

\begin{enumerate}
  	\item Démontrer que toute suite appartenant à $\mathscr{S}$ est strictement positive à partir du rang $1$.
  	\item Soit $\left(u_{n}\right)_{n \geq 0}$ une suite appartenant à $\mathscr{S}$. Démontrer que, s'il existe un rang $N \geq 2$ pour lequel $u_{N} \leq 1$, alors $\left(u_{n}\right)_{n \geq 0}$ converge vers 0.
  	\item Soit $\left(u_{n}\right)_{n \geq 0}$ une suite appartenant à $\mathscr{S}$. Démontrer que, si cette suite ne converge pas vers 0, alors elle diverge vers $+\infty$.
\end{enumerate}

\n Ci-dessous, on note $E_{0}$ l'ensemble des réels $x$ pour lesquels la suite $u(x)$ converge vers 0, et $E_{\infty}$ l'ensemble des réels $x$ pour lesquels $u(x)$ diverge vers $+\infty$.


\begin{enumerate}\setcounter{enumi}{3}
  	\item Démontrer que $0 \in E_{0}$.
  	\item \begin{enumerate}
  	\item Démontrer, pour tout entier $n \geq 0$, que la fonction $x \mapsto u_{n}(x)$ est strictement croissante sur $\mathbb{R}$.
  	\item En déduire que, si $x$ est un élément de $E_{0}$, alors l'intervalle $\left.]-\infty, x\right]$ est inclus dans $E_{0}$.
  	\end{enumerate}
  	\item \begin{enumerate}
  	\item Démontrer que la fonction $x \mapsto \exp (x)-x(x+1)$ est strictement positive sur l'intervalle $[2,+\infty[$.
  	\item Soit $\left(u_{n}\right)_{n \geq 0}$ une suite appartenant à $\mathscr{S}$. Démontrer que, s'il existe un rang $N \geq 1$ pour lequel $u_{N} \geq N+1$, alors $\left(u_{n}\right)_{n \geq 0}$ diverge vers $+\infty$.
  	\item Démontrer que $1 \in E_{\infty}$.
  	\end{enumerate}
	\item Démontrer que, si $x$ est un élément de $E_{\infty}$, alors l'intervalle $[x,+\infty[$ est inclus dans $E_{\infty}$.
\end{enumerate}

\n Nous allons maintenant démontrer qu'il existe un nombre réel $\delta$ tel que l'intervalle $] - \infty, \delta [$ est inclus dans $E_{0}$ et l'intervalle $[\delta,+\infty[$ est inclus dans $E_{\infty}$.

\begin{enumerate}\setcounter{enumi}{7}
  	\item On définit deux suites $\left(a_{n}\right)_{n \geq 0}$ et $\left(b_{n}\right)_{n \geq 0}$ de la façon suivante.\\
  	Tout d'abord, on pose $a_{0}=0$ et $b_{0}=1$.\\
  	Puis, pour tout entier $n \geq 0$, on pose $a_{n+1}=\left(a_{n}+b_{n}\right) / 2$ et $b_{n+1}=b_{n}$ si $\left(a_{n}+b_{n}\right) / 2 \in E_{0}$, et on pose $a_{n+1}=a_{n}$ et $b_{n+1}=\left(a_{n}+b_{n}\right) / 2$ sinon.
  	\begin{enumerate}
  		\item Démontrer que les suites $\left(a_{n}\right)_{n \geq 0}$ et $\left(b_{n}\right)_{n \geq 0}$ sont convergentes et ont même limite.
  		\item Soit $\delta$ la limite commune aux suites $\left(a_{n}\right)_{n \geq 0}$ et $\left(b_{n}\right)_{n \geq 0}$. Démontrer que l'intervalle $]-\infty, \delta[$ est inclus dans $E_{0}$ et l'intervalle $] \delta,+\infty$ [ est inclus dans $E_{\infty}$.
  	\end{enumerate}
  	\item On pose $c_{2}=\ln (\ln (2)), c_{3}=\ln (\ln (2 \ln (3)))$ et $c_{4}=\ln (\ln (2 \ln (3 \ln (4))))$, et plus généralement, pour tout entier $\ell \geq 2, c_{\ell}=\ln (\ln (2 \ln (3 \ln (\cdots \ln ((\ell-1) \ln (\ell)) \ldots)))$.\\
  	Démontrer que, pour tout entier $\ell \geq 2$, le nombre réel $c_{\ell}$ appartient à $E_{0}$.
  	\item Démontrer que la suite $\left(c_{\ell}\right)_{\ell \geq 2}$ converge.
  	\item Démontrer que $\delta \in E_{\infty}$.
\end{enumerate}


\section{Problème 2 : La loi du milieu}

\n Soit $n$ un entier naturel non nul. Dans un sac, on place $2 n+1$ boules indiscernables au toucher et numérotées $0,1,2, \ldots, 2 n$. On vide alors progressivement le sac jusqu'à n'y laisser qu'une seule boule, selon le protocole suivant :

\begin{itemize}
  \item on tire trois boules simultanément;
  \item si les trois boules tirées ont pour numéros $a, b$ et $c$, avec $a<b<c$, on élimine les boules de numéros $a$ et $c$ et on replace dans le sac la boule de numéro $b$;
  \item on recommence les opérations précédentes.
\end{itemize}

\n Au bout de $n$ tirages, il ne reste plus qu'une seule boule, et on note $D_{n}$ son numéro. Pour tout entier $k$, on note $\Prob\left[D_{n}=k\right]$ la probabilité que la dernière boule restant dans le sac soit celle de numéro $k$.

\subsection{Étude des petits cas}

\begin{enumerate}
  	\item Déterminer la loi de la variable aléatoire $D_{1}$.
  	\item Déterminer la loi de la variable aléatoire $D_{2}$.
\end{enumerate}

\subsection{Valeurs extrêmes et symétrie}


\begin{enumerate}\setcounter{enumi}{2}
  	\item Déterminer la probabilité $\Prob\left[D_{n}=0\right]$.
	\item Déterminer la probabilité $\Prob\left[D_{n}=1\right]$ en fonction de $n$.
  	\item Soit $i$ un entier tel que $0 \leq i \leq 2 n$. Pourquoi a-t-on $\Prob\left[D_{n}=i\right]=\Prob\left[D_{n}=2 n-i\right]$ ?
  \item Calculer l'espérance de la variable aléatoire $D_{n}$ en fonction de $n$.
\end{enumerate}

\subsection{Comportement limite}

\n Dans cette partie, on souhaite étudier la loi de $D_{n}$ lorsque $n$ tend vers $+\infty$. Afin de faciliter cette étude, on démontre tout d'abord un résultat préliminaire.


\begin{enumerate}\setcounter{enumi}{6}
  	\item On considère la suite $\left(u_{n}\right)_{n \geq 0}$ définie par $u_{0}=1$ et par
  	$$u_{n}=\frac{1}{2} \times \frac{3}{4} \times \frac{5}{6} \times \cdots \times \frac{2 n-1}{2 n}$$
  	pour tout $n \geq 1$. Démontrer que $u_{n} \leq \frac{1}{\sqrt{3 n+1}}$ pour tout $n \geq 0$.
\end{enumerate}

\n Il est maintenant temps d'étudier la loi de $D_{n}$ elle-même.

\begin{enumerate}\setcounter{enumi}{7}
  	\item Déterminer, pour tout entier $j$ tel que $0 \leq j \leq 2 n$, la probabilité $p_{j}$ que la boule de numéro $j$ soit éliminée lors de la première sélection.
  	\item Démontrer que, si $n \geq 3$, alors $p_{j} \geq \frac{1}{2 n}$ pour tout entier $j$ tel que $0 \leq j \leq 2 n$.
  	\item On note $M_{n}$ la plus grande des probabilités $\Prob\left[D_{n}=j\right]$ lorsque $0 \leq j \leq 2 n$.\\
  	Démontrer que $M_{n}$ tend vers 0 lorsque $n$ tend vers $+\infty$.
\end{enumerate}

\subsection{Résultat le plus probable}

\n On rappelle que, pour deux événements $A$ et $B$, on note $A \setminus B$ l'événement selon lequel $A$ est réalisé, mais pas $B$. En outre, si $\Prob[B] \neq 0$, on note $\Prob_{B}[A]$ la probabilité conditionnelle de $A$ sachant $B$.\\

\n On souhaite ici démontrer, pour tout entier $n \geq 1$, que $\Prob\left[D_{n}=n\right]=M_{n}$. Dans ce but, on va démontrer la propriété $\mathscr{P}_{n}$ suivante :

\begin{center}
Pour tout entier $k$ tel que $0 \leq k \leq n-1$, on a $\Prob\left[D_{n}=k\right] \leq \Prob\left[D_{n}=k+1\right]$.
\end{center}


\begin{enumerate}\setcounter{enumi}{10}
  	\item Démontrer que, si $\mathscr{P}_{n}$ est vraie, alors $\Prob\left[D_{n}=n\right]=M_{n}$.
  	\item Démontrer $\mathscr{P}_{1}$.
\end{enumerate}

\n On suppose maintenant que l'on dispose d'un entier $n \geq 2$ tel que $\mathscr{P}_{n-1}$ est vraie et d'un entier $k$ tel que $0 \leq k \leq n-1$.


\begin{enumerate}\setcounter{enumi}{12}
	\item Pour tout entier $\ell$ compris entre 0 et $2 n$, distinct de $k$ et de $k+1$, on note $X_{\ell}$ l'événement selon lequel les trois boules de numéros $k, k+1$ et $\ell$ sont choisies dès la première sélection.
	\begin{enumerate}
		\item Pourquoi, si $\ell>k+1$, a-t-on $\Prob_{X_{\ell}}\left[D_{n}=k\right]=0$ et $\Prob_{X_{\ell}}\left[D_{n}=k+1\right]=\Prob\left[D_{n-1}=k\right]$ ?
		\item Donner des résultats analogues sur $\Prob_{X_{\ell}}\left[D_{n}=k\right]$ et $\Prob_{X_{\ell}}\left[D_{n}=k+1\right]$ lorsque $\ell<k$.
		\item On note maintenant $X$ l'événement selon lequel les deux boules de numéros $k$ et $k+1$ sont choisies dès la première sélection. Démontrer que $\Prob_{X}\left[D_{n}=k\right] \leq \Prob_{X}\left[D_{n}=k+1\right]$.
	\end{enumerate}
	\item Soit $Y$ l'événement selon lequel l'une des boules de numéros $k$ et $k+1$ est éliminée lors de la première sélection.
	\begin{enumerate}
		\item Démontrer que $\Prob_{Y \setminus X}\left[D_{n}=k\right]=\Prob_{Y \setminus X}\left[D_{n}=k+1\right]$.
		\item En déduire que $\Prob_{Y}\left[D_{n}=k\right] \leq \Prob_{Y}\left[D_{n}=k+1\right]$.
	\end{enumerate}
	\item Soient $a, b$ et $c$ les numéros des trois boules choisies lors de la première sélection, avec $a<b<c$.
	\begin{enumerate}
		\item Soit $G$ l'événement selon lequel $c<k$.\\ 
		Démontrer que $\Prob_{G}\left[D_{n}=k\right] \leq \Prob_{G}\left[D_{n}=k+1\right]$.
		\item Soit $H$ l'événement selon lequel $a<k$ et $k+1<c$.\\ 
		Démontrer que $\Prob_{H}\left[D_{n}=k\right] \leq \Prob_{H}\left[D_{n}=k+1\right]$.
		\item Soit $I$ l'événement selon lequel $k+1<a$. Démontrer que, si $k \leq n-2$, alors :
		$$\Prob_{I}\left[D_{n}=k\right] \leq \Prob_{I}\left[D_{n}=k+1\right]$$
	\end{enumerate}
	\item Démontrer que, si $k \leq n-2$, alors $\Prob\left[D_{n}=k\right] \leq \Prob\left[D_{n}=k+1\right]$.
	\item Démontrer $\mathscr{P}_{n}$.
\end{enumerate}


\section{Problème 3 : Que la force soit avec $f$ !}

\n Dans tout le problème, $k$ désigne un entier naturel non nul, $I$ un intervalle ouvert de $] 0,+\infty[$, et $f$ une fonction définie sur $I$ et à valeurs strictement positives.\\
On dit que la fonction $f$ est << $k$-forte >> si, pour tous les réels $x$ et $y$ appartenant à $I$,
$$\left(y^{k} f(y)-x^{k} f(x)\right)\left(\frac{f(y)}{y^{k}}-\frac{f(x)}{x^{k}}\right) \geq 0$$
On dit que $f$ est << $k$-faible >> si, pour tous les réels $x$ et $y$ appartenant à $I$,
$$\left(y^{k} f(y)-x^{k} f(x)\right)\left(\frac{f(y)}{y^{k}}-\frac{f(x)}{x^{k}}\right) \leq 0$$

\subsection{Quelques exemples et propriétés}

\begin{enumerate}
	\item Démontrer que la fonction $f_{1}$ définie sur l'intervalle $]0,+\infty[$ par $f_{1}(x)=x^{2}$ est 1-forte et 3-faible.
	\item Démontrer que la fonction $f_{2}$ définie sur l'intervalle $]0,1[$ par $f_{2}(x)=\exp(x)$ est 1-faible mais pas 1-forte.
	\item Démontrer que la fonction $f_{3}$ définie sur l'intervalle $]1,+\infty[$ par $f_{3}(x)=\exp(x)$ est 1-forte mais pas 1-faible.
	\item Démontrer que la fonction $f_{4}$ définie sur l'intervalle $ 0,+\infty[$ par $f_{4}(x)=\frac{1}{x}$ est $k$-faible pour tout entier $k \geq 1$.
	\item Existe-t-il une fonction définie sur l'intervalle $]0,+\infty[$ qui soit $k$-forte pour tout entier $k \geq 1$ ?
\end{enumerate}

\subsection{Quelques critères de force et de faiblesse}


\begin{enumerate}\setcounter{enumi}{5}
	\item Démontrer que $f$ est $k$-forte si et seulement si
	$$\frac{x^{k}}{y^{k}}+\frac{y^{k}}{x^{k}} \leq \frac{f(x)}{f(y)}+\frac{f(y)}{f(x)}$$
	pour tous les réels $x$ et $y$ appartenant à $I$, et que $f$ est $k$-faible si et seulement si
	$$\frac{x^{k}}{y^{k}}+\frac{y^{k}}{x^{k}} \geq \frac{f(x)}{f(y)}+\frac{f(y)}{f(x)}$$
	pour tous les réels $x$ et $y$ appartenant à $I$.
	\item Démontrer que $f$ est $k$-forte si et seulement si
	$$\frac{\max \left(x^{k}, y^{k}\right)}{\min \left(x^{k}, y^{k}\right)} \leq \frac{\max (f(x), f(y))}{\min (f(x), f(y))}$$
	pour tous les réels $x$ et $y$ appartenant à $I$, et que $f$ est $k$-faible si et seulement si
	$$\frac{\max \left(x^{k}, y^{k}\right)}{\min \left(x^{k}, y^{k}\right)} \geq \frac{\max (f(x), f(y))}{\min (f(x), f(y))}$$
	pour tous les réels $x$ et $y$ appartenant à $I$.
	\item On note $g_{k}$ et $h_{k}$ les fonctions définies sur $I$ par :
	$$g_{k}(x)=x^{k} f(x) \q \text { et }\q h_{k}(x)=\frac{f(x)}{x^{k}} $$
	\begin{enumerate}
		\item Démontrer que, si $g_{k}$ et $h_{k}$ sont monotones, alors $f$ est $k$-forte ou $k$-faible.
		\item Démontrer que, si $f$ est $k$-faible, alors $g_{k}$ et $h_{k}$ sont monotones.
		\item Soit $f$ la fonction définie sur l'intervalle $I=] 0,+\infty[$ par
	$$f(x)=  \left\{\begin{array}{ll} 
        x & \text { si } 0<x<1\\ 
        4 & \text { si } x=1\\
        x & \text { si } 1<x<2\\
        4x & \text { si } 2 \leq x
      \end{array}  
    \right.$$
	Démontrer que $f$ est 1-forte mais que les fonctions $g_{1}$ et $h_{1}$ ne sont pas monotones.
	\end{enumerate}
	\item On suppose dans cette question que $f$ est dérivable sur $I$ et que sa dérivée $f^{\prime}$ est continue sur $I$.
	\begin{enumerate}
		\item Démontrer que, si $\left|f^{\prime}(x)\right| \geq k \frac{f(x)}{x}$ pour tout réel $x \in I$, alors $f$ est $k$-forte.
		\item Démontrer que, si $\left|f^{\prime}(x)\right| \leq k \frac{f(x)}{x}$ pour tout réel $x \in I$, alors $f$ est $k$-faible.
		\item Démontrer que les réciproques aux questions \textbf{9.a.} et \textbf{9.b.} sont vraies.
	\end{enumerate}
\end{enumerate}


\subsection{Une multitude de fonctions fortes et faibles}

\n On dit que la fonction $f$ est << forte >> s'il existe un entier $k \geq 1$ pour lequel $f$ est $k$-forte, et que $f$ est << faible >> s'il existe un entier $k \geq 1$ pour lequel $f$ est $k$-faible.


\begin{enumerate}\setcounter{enumi}{9}
	\item Démontrer que, si $f$ est faible, la fonction $F$ définie sur $I$ par $F(x)=\frac{1}{f(x)}$ est faible.
	\item Démontrer que, si deux fonctions $f$ et $g$ définies sur $I$ sont faibles, les fonctions $f+g$, $f \times g$ et $\frac{f}{g}$ sont faibles.
	\item Démontrer à l'aide de contre-exemples que, si deux fonctions $f$ et $g$ définies sur $I$ sont fortes, les fonctions $f+g, f \times g$ et $\frac{f}{g}$ ne sont pas nécessairement fortes.
	\item Soit $f$ une fonction définie sur $I$ et à valeurs strictement positives, et $g$ une fonction définie sur $]0,+\infty[$.
	\begin{enumerate}
		\item Démontrer que, si $f$ et $g$ sont faibles, la fonction $g \circ f$ est faible.
		\item Démontrer que, si $f$ et $g$ sont fortes, la fonction $g \circ f$ est forte.
	\end{enumerate}
\end{enumerate}


\subsection{Application à la démonstration d'inégalités}


\begin{enumerate}\setcounter{enumi}{13}
	\item Soit $a, b$ et $c$ trois réels strictement positifs, et $n$ un entier naturel non nul. Démontrer que :
	$$\left(\frac{a+c}{b+c}\right)^{n}+\left(\frac{b+c}{a+c}\right)^{n} \leq\left(\frac{a}{b}\right)^{n}+\left(\frac{b}{a}\right)^{n}$$
	\item Dans cette question, on pourra utiliser le fait que les fonctions cos et sin sont dérivables sur $] 0, \frac{\pi}{2}[$ de dérivées respectivement $\cos'=-\sin$ et $\sin'=\cos$.\\
	La fonction tan est définie sur $\left]0,\frac{\pi}{2}\right[$ par $\tan (x)=\frac{\sin (x)}{\cos (x)}$.\\
	Soient $a$ et $b$ deux nombres réels de l'intervalle $]0, \frac{\pi}{2}[$, démontrer que :
	$$\frac{\sin (a)}{\sin (b)}+\frac{\sin (b)}{\sin (a)} \leq \frac{a}{b}+\frac{b}{a} \leq \frac{\tan (a)}{\tan (b)}+\frac{\tan (b)}{\tan (a)}$$
\end{enumerate}

\end{document}